\documentclass[11pt,a4paper]{article}

\usepackage[utf8]{inputenc}
\usepackage[T1]{fontenc}
\usepackage[french]{babel}
\usepackage{amsmath,amssymb}
\usepackage{geometry}
\usepackage{enumitem}
\usepackage{multicol}
\usepackage{tabularx}
\usepackage{booktabs}
\usepackage{xcolor}
\usepackage{mdframed}
\usepackage{lmodern}
\usepackage{microtype}
\usepackage{fancyhdr}
\usepackage{lastpage}
\usepackage{parskip}

\geometry{top=2cm, bottom=2.5cm, left=2.2cm, right=2.2cm}

\pagestyle{fancy}
\fancyhf{}
\fancyhead[L]{\small\textit{Mathématiques - Corrigé pédagogique}}
\fancyhead[R]{\small\textit{Épreuve de synthèse - Analyse}}
\fancyfoot[C]{\small Page \thepage\ / \pageref{LastPage}}
\renewcommand{\headrulewidth}{0.3pt}
\renewcommand{\footrulewidth}{0pt}

% Couleurs
\definecolor{grisbox}{gray}{0.94}
\definecolor{bleutexte}{RGB}{30,60,120}
\definecolor{vertok}{RGB}{0,120,60}
\definecolor{rougeerr}{RGB}{180,30,30}
\definecolor{orangeinfo}{RGB}{200,100,0}
\definecolor{bleuclair}{RGB}{220,230,245}
\definecolor{vertclair}{RGB}{220,245,230}
\definecolor{orangeclair}{RGB}{255,240,220}

% Environnements colorés
\newmdenv[backgroundcolor=vertclair, linecolor=vertok, linewidth=1pt,
          innerleftmargin=10pt, innerrightmargin=10pt,
          innertopmargin=6pt, innerbottommargin=6pt,
          skipabove=4pt, skipbelow=6pt]{reponsebox}

\newmdenv[backgroundcolor=bleuclair, linecolor=bleutexte, linewidth=1pt,
          innerleftmargin=10pt, innerrightmargin=10pt,
          innertopmargin=6pt, innerbottommargin=6pt,
          skipabove=4pt, skipbelow=4pt]{methodebox}

\newmdenv[backgroundcolor=orangeclair, linecolor=orangeinfo, linewidth=1pt,
          innerleftmargin=10pt, innerrightmargin=10pt,
          innertopmargin=6pt, innerbottommargin=6pt,
          skipabove=4pt, skipbelow=6pt]{piegebox}

\newmdenv[backgroundcolor=grisbox, linewidth=0pt,
          innerleftmargin=10pt, innerrightmargin=10pt,
          innertopmargin=6pt, innerbottommargin=6pt,
          skipabove=6pt, skipbelow=6pt]{notebox}

% Style des sections
\usepackage{titlesec}
\titleformat{\section}[block]{\large\bfseries\color{bleutexte}}{}{0pt}{}[\vspace{-4pt}\rule{\linewidth}{0.4pt}]
\titlespacing{\section}{0pt}{16pt}{8pt}

% En-tête de question
\newcounter{qnum}
\newcommand{\qtitre}[2]{%
  \stepcounter{qnum}%
  \medskip\noindent
  \textbf{\color{bleutexte}Question \theqnum}\enspace
  \textit{(#1)}\enspace
  \textcolor{vertok}{\textbf{$\Rightarrow$ Réponse : #2}}
  \par\smallskip
}

% Énoncé de la question
\newcommand{\qenonce}[1]{%
  \begin{mdframed}[backgroundcolor=grisbox, linewidth=0pt,
    innerleftmargin=8pt, innerrightmargin=8pt,
    innertopmargin=5pt, innerbottommargin=5pt,
    skipabove=2pt, skipbelow=4pt]
  \small\textit{#1}
  \end{mdframed}
}

\begin{document}

% ─── TITRE ───────────────────────────────────────────────────────────────────
\begin{center}
  {\Large\bfseries\color{bleutexte} Corrigé - Épreuve de synthèse Mathématiques}\\[4pt]
  {\large Fonctions de plusieurs variables, Opérateurs vectoriels \& Équations différentielles}\\[6pt]
  {\small\itshape
    Pour chaque question : la réponse correcte, la méthode complète et les pièges à éviter.}
\end{center}

\vspace{2pt}
\noindent\rule{\linewidth}{0.6pt}
\vspace{4pt}

% ── Grille de correction ──────────────────────────────────────────────────────
\begin{notebox}
\textbf{Grille de correction rapide}\\[4pt]
\begin{tabular}{@{}lll@{\quad}lll@{\quad}lll@{}}
Q1  & QCM    & \textbf{A}                             &
Q11 & QCM    & \textbf{A}                             &
Q21 & QCM    & \textbf{C}                             \\
Q2  & QCM    & \textbf{D}                             &
Q12 & Calcul & $3x^2\sin y\;;\;x^3\cos y\;;\;3x^2\cos y$ &
Q22 & QCM    & \textbf{A}                             \\
Q3  & QCM    & \textbf{C}                             &
Q13 & QCM    & \textbf{C}                             &
Q23 & QCM    & \textbf{B}                             \\
Q4  & QCM    & \textbf{B}                             &
Q14 & QCM    & \textbf{A}                             &
Q24 & Calcul & $r_1=3,\,r_2=-2$                       \\
Q5  & QCM    & \textbf{D}                             &
Q15 & QCM    & \textbf{B}                             &
Q25 & QCM    & \textbf{D}                             \\
Q6  & Calcul & $2xe^y\;;\;x^2e^y$                     &
Q16 & QCM    & \textbf{D}                             &
Q26 & QCM    & \textbf{B}                             \\
Q7  & QCM    & \textbf{A}                             &
Q17 & QCM    & \textbf{A}                             &
Q27 & QCM    & \textbf{C}                             \\
Q8  & QCM    & \textbf{C}                             &
Q18 & Calcul & $3x+z$                                 &
Q28 & QCM    & \textbf{D}                             \\
Q9  & QCM    & \textbf{B}                             &
Q19 & QCM    & \textbf{C}                             &
Q29 & QCM    & \textbf{C}                             \\
Q10 & QCM    & \textbf{D}                             &
Q20 & QCM    & \textbf{D}                             &
Q30 & Calcul & $8$                                    \\
\end{tabular}
\end{notebox}

% ═══════════════════════════════════════════════════════════════════════════════
\section{Partie 1 - Fonctions de plusieurs variables - définitions fondamentales}
% ═══════════════════════════════════════════════════════════════════════════════

%────── Q1 ────────────────────────────────────────────────────────────────────
\qtitre{QCM}{A - $\lim_{(x,y)\to a} f(x,y) = f(a)$}
\qenonce{Soit $f:\mathbb{R}^2\to\mathbb{R}$. Elle est continue en $a=(a_1,a_2)$ si et seulement si :
\quad A)~$\lim_{(x,y)\to a}f(x,y) = f(a)$
\quad B)~Elle admet des dérivées partielles en $a$
\quad C)~Elle est bornée au voisinage de $a$
\quad D)~$f(a) = 0$}

\begin{methodebox}
\textbf{Rappel.} La continuité en $a$ est définie par l'égalité entre la limite et la valeur :
$$f \text{ continue en } a \;\Longleftrightarrow\; \lim_{(x,y)\to a} f(x,y) = f(a)$$
\end{methodebox}

\begin{itemize}[leftmargin=1.5em, itemsep=2pt]
  \item \textcolor{vertok}{\textbf{A)}} Définition exacte de la continuité. \textbf{Correct.}
  \item \textcolor{rougeerr}{\textbf{B)}} L'existence des dérivées partielles ne suffit \emph{pas} à garantir la continuité (contre-exemple : $f(x,y)=\frac{xy}{x^2+y^2}$).
  \item \textcolor{rougeerr}{\textbf{C)}} Être borné n'implique pas la continuité.
  \item \textcolor{rougeerr}{\textbf{D)}} Annuler $f$ en $a$ n'a rien à voir avec la continuité.
\end{itemize}

%────── Q2 ────────────────────────────────────────────────────────────────────
\qtitre{QCM}{D - $F(x,y,z) = 0$}
\qenonce{Une surface de $\mathbb{R}^3$ donnée sous forme implicite s'écrit :
\quad A)~$\mathbf{r}(t) = (x(t),y(t),z(t))$
\quad B)~$z = f(x,y)$
\quad C)~$(u,v)\mapsto(r_1(u,v),r_2(u,v),r_3(u,v))$
\quad D)~$F(x,y,z) = 0$}

\begin{methodebox}
\textbf{Trois représentations d'une surface :}
\begin{itemize}[itemsep=1pt, topsep=2pt]
  \item \textbf{Explicite :} $z = f(x,y)$ - la surface est le graphe de $f$.
  \item \textbf{Implicite :} $F(x,y,z) = 0$ - ensemble de niveau d'une fonction de 3 variables.
  \item \textbf{Paramétrique :} $(u,v)\mapsto\mathbf{r}(u,v)$ - deux paramètres décrivent la surface.
\end{itemize}
La forme A, $\mathbf{r}(t)$, est un paramétrage de \emph{courbe} (une seule variable).
\end{methodebox}

%────── Q3 ────────────────────────────────────────────────────────────────────
\qtitre{QCM}{C - La pente de la tangente à $x\mapsto f(x,b)$ en $a$, $y$ fixé à $b$}
\qenonce{La dérivée partielle $\dfrac{\partial f}{\partial x}(a,b)$ représente :
\quad A)~La différentielle totale de $f$ en $(a,b)$
\quad B)~Le taux de variation de $f$ dans toutes les directions en $(a,b)$
\quad C)~La pente de la tangente à la courbe $x\mapsto f(x,b)$ en $a$, $y$ fixé à $b$
\quad D)~La pente de la tangente à la courbe $y\mapsto f(a,y)$ en $b$}

\begin{methodebox}
\textbf{Dérivée partielle.} $\dfrac{\partial f}{\partial x}(a,b)$ est la dérivée ordinaire de la fonction d'une variable $x\mapsto f(x,b)$, calculée en $x=a$, $y$ étant maintenu fixé à $b$.
\end{methodebox}

\begin{itemize}[leftmargin=1.5em, itemsep=2pt]
  \item \textcolor{rougeerr}{\textbf{A)}} La différentielle totale $df$ fait intervenir les deux dérivées partielles.
  \item \textcolor{rougeerr}{\textbf{B)}} Le taux de variation dans \emph{toutes les directions} est la dérivée directionnelle, pas $\partial_x f$.
  \item \textcolor{vertok}{\textbf{C)}} Définition correcte - $y$ est fixé, seul $x$ varie. \textbf{Correct.}
  \item \textcolor{rougeerr}{\textbf{D)}} Cette description correspond à $\dfrac{\partial f}{\partial y}(a,b)$.
\end{itemize}

%────── Q4 ────────────────────────────────────────────────────────────────────
\qtitre{QCM}{B - Continue en $(a,b)$}
\qenonce{Si $f$ est différentiable en $(a,b)$, alors $f$ est nécessairement :
\quad A)~De classe $\mathcal{C}^1$ au voisinage de $(a,b)$
\quad B)~Continue en $(a,b)$
\quad C)~Deux fois dérivable en $(a,b)$
\quad D)~Analytique en $(a,b)$}

\begin{methodebox}
\textbf{Hiérarchie des régularités (du plus fort au plus faible) :}
$$\mathcal{C}^\infty \;\Rightarrow\; \mathcal{C}^2 \;\Rightarrow\; \mathcal{C}^1 \;\Rightarrow\; \text{différentiable} \;\Rightarrow\; \text{continue} \;\Rightarrow\; \text{dérivées partielles existent}$$
La différentiabilité implique la continuité, mais pas la réciproque.
\end{methodebox}

\begin{itemize}[leftmargin=1.5em, itemsep=2pt]
  \item \textcolor{rougeerr}{\textbf{A)}} $\mathcal{C}^1$ implique différentiabilité, mais pas l'inverse.
  \item \textcolor{vertok}{\textbf{B)}} Différentiable $\Rightarrow$ continue. \textbf{Correct.}
  \item \textcolor{rougeerr}{\textbf{C)}} La différentiabilité ne garantit pas la dérivabilité au 2\textsuperscript{e} ordre.
  \item \textcolor{rougeerr}{\textbf{D)}} L'analyticité est bien plus forte (développement en série entière).
\end{itemize}

%────── Q5 ────────────────────────────────────────────────────────────────────
\qtitre{QCM}{D - Trouver deux chemins aboutissant à des limites différentes}
\qenonce{Pour montrer que $\lim_{(x,y)\to(0,0)} f(x,y)$ n'existe pas, on peut :
\quad A)~Vérifier que $f(0,0)$ n'est pas définie
\quad B)~Montrer que $f$ est non bornée au voisinage de $(0,0)$
\quad C)~Calculer les dérivées partielles en $(0,0)$
\quad D)~Trouver deux chemins aboutissant à des limites différentes}

\begin{methodebox}
\textbf{Méthode des chemins.} Si deux chemins vers $(0,0)$ donnent des limites différentes, la limite globale n'existe pas. L'existence ou la non-existence de $f(0,0)$, ou la non-bornitude, ne permettent pas à elles seules de conclure sur la limite.
\end{methodebox}

\begin{piegebox}
\textbf{Piège.} Tester un seul chemin (ex. $y=0$) peut donner une valeur finie, mais d'autres chemins peuvent mener à une valeur différente. La non-existence se prouve par deux chemins discordants ; l'existence nécessite une preuve globale (passage en polaires, encadrement...).
\end{piegebox}

%────── Q6 (Calcul) ───────────────────────────────────────────────────────────
\qtitre{Calcul}{$\dfrac{\partial f}{\partial x} = 2xe^y$ \quad et \quad $\dfrac{\partial f}{\partial y} = x^2e^y$}
\qenonce{Soit $f(x,y) = x^2 e^{y}$. Calculer $\dfrac{\partial f}{\partial x}$ et $\dfrac{\partial f}{\partial y}$.}

\begin{reponsebox}
On traite l'autre variable comme une constante à chaque fois.

\textbf{Dérivée par rapport à $x$} ($y$ constant) :
$$\boxed{\frac{\partial f}{\partial x} = 2x\,e^y}$$

\textbf{Dérivée par rapport à $y$} ($x$ constant) :
$$\boxed{\frac{\partial f}{\partial y} = x^2 e^y}$$

\textbf{Remarque :} ici $\dfrac{\partial f}{\partial y} = f$ car $e^y$ est sa propre dérivée par rapport à $y$.
\end{reponsebox}

% ═══════════════════════════════════════════════════════════════════════════════
\section{Partie 2 - Différentielle, règle des chaînes et plan tangent}
% ═══════════════════════════════════════════════════════════════════════════════

%────── Q7 ────────────────────────────────────────────────────────────────────
\qtitre{QCM}{A - $df = 2xy\,dx + (x^2+\cos y)\,dy$}
\qenonce{La différentielle de $f(x,y) = x^2 y + \sin(y)$ est :
\quad A)~$df = 2xy\,dx + (x^2+\cos y)\,dy$
\quad B)~$df = 2y\,dx + x^2\,dy$
\quad C)~$df = (2xy + x^2)\,d(x+y)$
\quad D)~$df = 2xy\,dx + x^2\,dy$}

\begin{methodebox}
\textbf{Différentielle.} $df = \dfrac{\partial f}{\partial x}dx + \dfrac{\partial f}{\partial y}dy$
\end{methodebox}

Calcul des dérivées partielles :
$$\frac{\partial f}{\partial x} = 2xy \qquad (\sin y \text{ est une constante par rapport à } x)$$
$$\frac{\partial f}{\partial y} = x^2 + \cos y \qquad (\text{dérivée de } x^2y \text{ puis de } \sin y \text{ par rapport à } y)$$

\textbf{Pourquoi D est faux :} oubli du terme $\cos y$ provenant de $\dfrac{\partial(\sin y)}{\partial y} = \cos y$.

%────── Q8 ────────────────────────────────────────────────────────────────────
\qtitre{QCM}{C - $\dfrac{\partial f}{\partial x} = g'(h)\cdot\dfrac{\partial h}{\partial x}$}
\qenonce{Si $f(x,y) = g(h(x,y))$ (composition), la règle de la chaîne donne :
\quad A)~$\dfrac{\partial f}{\partial x} = g'(h) + \dfrac{\partial h}{\partial x}$
\quad B)~$\dfrac{\partial f}{\partial x} = g(h)\cdot h(x,y)$
\quad C)~$\dfrac{\partial f}{\partial x} = g'(h)\cdot\dfrac{\partial h}{\partial x}$
\quad D)~$\dfrac{\partial f}{\partial x} = g''(h)$}

\begin{methodebox}
\textbf{Règle de la chaîne.} Si $f = g \circ h$ avec $h : \mathbb{R}^2 \to \mathbb{R}$ et $g : \mathbb{R} \to \mathbb{R}$ :
$$\frac{\partial f}{\partial x} = g'(h(x,y))\cdot\frac{\partial h}{\partial x}, \qquad \frac{\partial f}{\partial y} = g'(h(x,y))\cdot\frac{\partial h}{\partial y}$$
On dérive la fonction extérieure $g$ évaluée en $h$, puis on multiplie par la dérivée de $h$.
\end{methodebox}

\textbf{Exemple :} $f(x,y) = \sin(x^2 y)$, $h = x^2 y$, $g = \sin$ :
$$\frac{\partial f}{\partial x} = \cos(x^2y)\cdot 2xy$$

%────── Q9 ────────────────────────────────────────────────────────────────────
\qtitre{QCM}{B - $\dfrac{\partial^2 f}{\partial x\partial y} = \dfrac{\partial^2 f}{\partial y\partial x}$}
\qenonce{Le théorème de Schwartz affirme que, pour $f \in \mathcal{C}^2$ :
\quad A)~$\dfrac{\partial^2 f}{\partial x^2} = \dfrac{\partial^2 f}{\partial y^2}$
\quad B)~$\dfrac{\partial^2 f}{\partial x\partial y} = \dfrac{\partial^2 f}{\partial y\partial x}$
\quad C)~$\dfrac{\partial f}{\partial x} = \dfrac{\partial f}{\partial y}$
\quad D)~$\Delta f = 0$}

\begin{methodebox}
\textbf{Théorème de Clairaut-Schwartz.} Pour $f\in\mathcal{C}^2$ sur un ouvert, l'ordre de dérivation n'a pas d'importance :
$$\frac{\partial^2 f}{\partial x\partial y} = \frac{\partial^2 f}{\partial y\partial x}$$
\textbf{Conséquence pratique :} pour calculer une dérivée croisée, on peut dériver dans l'ordre le plus simple.
\end{methodebox}

\begin{itemize}[leftmargin=1.5em, itemsep=2pt]
  \item \textcolor{rougeerr}{\textbf{A)}} Faux en général : $f_{xx} = f_{yy}$ n'est pas une propriété générale.
  \item \textcolor{vertok}{\textbf{B)}} Exactement l'énoncé du théorème. \textbf{Correct.}
  \item \textcolor{rougeerr}{\textbf{C, D)}} Sans rapport avec le théorème de Schwartz.
\end{itemize}

%────── Q10 ────────────────────────────────────────────────────────────────────
\qtitre{QCM}{D - $z = f(a,b) + \dfrac{\partial f}{\partial x}(a,b)(x-a) + \dfrac{\partial f}{\partial y}(a,b)(y-b)$}
\qenonce{L'équation du plan tangent à la surface $z = f(x,y)$ au point $(a,b,f(a,b))$ est :
\quad A)~$z = f(a,b)\cdot(x+y)$
\quad B)~$z = \dfrac{\partial f}{\partial x}(a,b)\,x + \dfrac{\partial f}{\partial y}(a,b)\,y$
\quad C)~$z = f'(a)(x-a)$
\quad D)~$z = f(a,b) + \dfrac{\partial f}{\partial x}(a,b)(x-a) + \dfrac{\partial f}{\partial y}(a,b)(y-b)$}

\begin{methodebox}
\textbf{Plan tangent = développement limité d'ordre 1.} Le plan tangent en $(a,b)$ est l'approximation affine :
$$z = f(a,b) + \frac{\partial f}{\partial x}(a,b)(x-a) + \frac{\partial f}{\partial y}(a,b)(y-b)$$
\end{methodebox}

La proposition B oublie le terme $f(a,b)$ (valeur de $f$ au point de tangence). La proposition C est une formule en une seule variable.

%────── Q11 ────────────────────────────────────────────────────────────────────
\qtitre{QCM}{A - $\nabla F(M_0)$}
\qenonce{Pour une surface $F(x,y,z)=0$, le vecteur normal en un point $M_0$ est :
\quad A)~$\nabla F(M_0)$
\quad B)~$\text{rot}\,\nabla F(M_0)$
\quad C)~$\Delta F(M_0)$
\quad D)~$(F_x + F_y + F_z)(M_0)$}

\begin{methodebox}
\textbf{Vecteur normal à une surface implicite.} La surface $F(x,y,z)=0$ est une surface de niveau de $F$. Le gradient $\nabla F$ est perpendiculaire aux surfaces de niveau, donc $\nabla F(M_0)$ est le \textbf{vecteur normal} à la surface en $M_0$.

La proposition B est nulle par Schwartz ($\text{rot}(\nabla F) = \mathbf{0}$). La proposition D est un scalaire, pas un vecteur.
\end{methodebox}

%────── Q12 (Calcul) ──────────────────────────────────────────────────────────
\qtitre{Calcul}{$\dfrac{\partial f}{\partial x} = 3x^2\sin(y)$,\; $\dfrac{\partial f}{\partial y} = x^3\cos(y)$,\; $\dfrac{\partial^2 f}{\partial x\partial y} = 3x^2\cos(y)$}
\qenonce{Soit $f(x,y) = x^3\sin(y)$. Calculer $\dfrac{\partial f}{\partial x}$, $\dfrac{\partial f}{\partial y}$ et $\dfrac{\partial^2 f}{\partial x\partial y}$. Vérifier le théorème de Schwartz.}

\begin{reponsebox}
\textbf{Dérivées partielles d'ordre 1 :}
$$\frac{\partial f}{\partial x} = 3x^2\sin(y) \qquad \frac{\partial f}{\partial y} = x^3\cos(y)$$

\textbf{Dérivée croisée (deux façons) :}
$$\frac{\partial^2 f}{\partial x\partial y} = \frac{\partial}{\partial y}\bigl(3x^2\sin y\bigr) = \boxed{3x^2\cos(y)}$$
$$\frac{\partial^2 f}{\partial y\partial x} = \frac{\partial}{\partial x}\bigl(x^3\cos y\bigr) = 3x^2\cos(y)$$

\textbf{Vérification :} $\dfrac{\partial^2 f}{\partial x\partial y} = \dfrac{\partial^2 f}{\partial y\partial x} = 3x^2\cos(y)$ \quad\checkmark \quad (Schwartz : $f\in\mathcal{C}^2$)
\end{reponsebox}

% ═══════════════════════════════════════════════════════════════════════════════
\section{Partie 3 - Opérateurs vectoriels - gradient, divergence, rotationnel}
% ═══════════════════════════════════════════════════════════════════════════════

%────── Q13 ────────────────────────────────────────────────────────────────────
\qtitre{QCM}{C - $\nabla f = \left(\dfrac{\partial f}{\partial x},\ \dfrac{\partial f}{\partial y},\ \dfrac{\partial f}{\partial z}\right)$}
\qenonce{Le gradient d'un champ scalaire $f(x,y,z)$ est :
\quad A)~$\nabla f = \dfrac{\partial^2 f}{\partial x^2} + \dfrac{\partial^2 f}{\partial y^2} + \dfrac{\partial^2 f}{\partial z^2}$
\quad B)~$\nabla f = \dfrac{\partial f}{\partial x} + \dfrac{\partial f}{\partial y} + \dfrac{\partial f}{\partial z}$
\quad C)~$\nabla f = \left(\dfrac{\partial f}{\partial x},\ \dfrac{\partial f}{\partial y},\ \dfrac{\partial f}{\partial z}\right)$
\quad D)~$\nabla f = f\cdot\left(\dfrac{\partial}{\partial x} + \dfrac{\partial}{\partial y} + \dfrac{\partial}{\partial z}\right)$}

\begin{methodebox}
\textbf{Gradient.} $\nabla f = \mathrm{grad}\,f = \left(\dfrac{\partial f}{\partial x},\;\dfrac{\partial f}{\partial y},\;\dfrac{\partial f}{\partial z}\right)$ est un \textbf{vecteur}.

La proposition A est le laplacien $\Delta f$ (scalaire). La proposition B est une somme scalaire, pas un vecteur.
\end{methodebox}

%────── Q14 ────────────────────────────────────────────────────────────────────
\qtitre{QCM}{A - Perpendiculairement aux surfaces de niveau, dans le sens des valeurs croissantes}
\qenonce{En tout point, le vecteur gradient $\nabla f$ est orienté :
\quad A)~Perpendiculairement aux surfaces de niveau, dans le sens des valeurs croissantes
\quad B)~Parallèlement aux surfaces de niveau de $f$
\quad C)~Dans la direction de l'axe $Oz$
\quad D)~Vers le minimum de $f$}

\begin{methodebox}
\textbf{Interprétation géométrique du gradient.}
\begin{itemize}[itemsep=1pt, topsep=2pt]
  \item $\nabla f$ est \textbf{perpendiculaire} aux surfaces de niveau $\{f = c\}$.
  \item $\nabla f$ pointe dans la \textbf{direction de plus forte croissance} de $f$.
  \item $\|\nabla f\|$ mesure le taux de variation maximal.
  \item $\nabla f = \mathbf{0}$ aux points critiques (extrema, selles).
\end{itemize}
\end{methodebox}

%────── Q15 ────────────────────────────────────────────────────────────────────
\qtitre{QCM}{B - $\dfrac{\partial P}{\partial x} + \dfrac{\partial Q}{\partial y} + \dfrac{\partial R}{\partial z}$}
\qenonce{La divergence de $\mathbf{F} = (P, Q, R)$ est :
\quad A)~$\sqrt{P^2+Q^2+R^2}$
\quad B)~$\dfrac{\partial P}{\partial x} + \dfrac{\partial Q}{\partial y} + \dfrac{\partial R}{\partial z}$
\quad C)~$\dfrac{\partial P}{\partial y} + \dfrac{\partial Q}{\partial z} + \dfrac{\partial R}{\partial x}$
\quad D)~Le vecteur $\left(\dfrac{\partial Q}{\partial z}-\dfrac{\partial R}{\partial y},\ldots\right)$}

\begin{methodebox}
\textbf{Divergence.} $\text{div}\,\mathbf{F} = \nabla\cdot\mathbf{F} = \dfrac{\partial P}{\partial x}+\dfrac{\partial Q}{\partial y}+\dfrac{\partial R}{\partial z}$ est un \textbf{scalaire}.

\textbf{Interprétation physique :} mesure le taux d'expansion local du champ (flux sortant par unité de volume).
\end{methodebox}

La proposition A est la norme de $\mathbf{F}$. La proposition C est une permutation cyclique erronée. La proposition D décrit le \textbf{rotationnel} (vecteur).

%────── Q16 ────────────────────────────────────────────────────────────────────
\qtitre{QCM}{D - $\text{rot}\,\mathbf{F} = \mathbf{0}$}
\qenonce{Si $\mathbf{F} = \nabla f$ est un champ gradient (champ conservatif), alors :
\quad A)~$\text{div}\,\mathbf{F} = 0$
\quad B)~$\Delta f = 0$
\quad C)~$\mathbf{F} = \mathbf{0}$
\quad D)~$\text{rot}\,\mathbf{F} = \mathbf{0}$}

\begin{methodebox}
\textbf{Propriété fondamentale :} $\text{rot}(\nabla f) = \mathbf{0}$ pour tout $f\in\mathcal{C}^2$.

\textbf{Preuve (composante $z$) :}
$$\frac{\partial}{\partial x}\frac{\partial f}{\partial y} - \frac{\partial}{\partial y}\frac{\partial f}{\partial x} = \frac{\partial^2 f}{\partial x\partial y} - \frac{\partial^2 f}{\partial y\partial x} = 0 \quad\text{(Schwartz)}$$

Un champ gradient est dit \textbf{irrotationnel} ou \textbf{conservatif}. La proposition A serait vraie pour un champ solénoïdal ($\text{div}=0$).
\end{methodebox}

%────── Q17 ────────────────────────────────────────────────────────────────────
\qtitre{QCM}{A - $\Delta f = \dfrac{\partial^2 f}{\partial x^2} + \dfrac{\partial^2 f}{\partial y^2} + \dfrac{\partial^2 f}{\partial z^2}$}
\qenonce{Le laplacien d'un champ scalaire $f(x,y,z)$ est défini par :
\quad A)~$\Delta f = \dfrac{\partial^2 f}{\partial x^2} + \dfrac{\partial^2 f}{\partial y^2} + \dfrac{\partial^2 f}{\partial z^2}$
\quad B)~$\Delta f = \left(\dfrac{\partial f}{\partial x}\right)^2 + \left(\dfrac{\partial f}{\partial y}\right)^2 + \left(\dfrac{\partial f}{\partial z}\right)^2$
\quad C)~$\Delta f = \text{rot}(\nabla f)$
\quad D)~$\Delta f = \nabla \cdot f$}

\begin{methodebox}
\textbf{Laplacien.} $\Delta f = \text{div}(\nabla f) = \dfrac{\partial^2 f}{\partial x^2}+\dfrac{\partial^2 f}{\partial y^2}+\dfrac{\partial^2 f}{\partial z^2}$ est un \textbf{scalaire}.

La proposition B confond les dérivées premières (au carré) et les dérivées secondes. La proposition C est nulle par Schwartz ($\text{rot}(\nabla f) = \mathbf{0}$).
\end{methodebox}

%────── Q18 (Calcul) ──────────────────────────────────────────────────────────
\qtitre{Calcul}{$\text{div}\,\mathbf{F} = 3x + z$}
\qenonce{Soit $\mathbf{F}(x,y,z) = (x^2,\ yz,\ xz)$. Calculer $\text{div}\,\mathbf{F}$.}

\begin{reponsebox}
Avec $P = x^2$, $Q = yz$, $R = xz$ :
$$\text{div}\,\mathbf{F} = \frac{\partial(x^2)}{\partial x} + \frac{\partial(yz)}{\partial y} + \frac{\partial(xz)}{\partial z} = 2x + z + x = \boxed{3x + z}$$
\end{reponsebox}

% ═══════════════════════════════════════════════════════════════════════════════
\section{Partie 4 - Équations différentielles - rappels de cours}
% ═══════════════════════════════════════════════════════════════════════════════

%────── Q19 ────────────────────────────────────────────────────────────────────
\qtitre{QCM}{C - $y' + 3y = e^x$}
\qenonce{Laquelle de ces équations est une EDO linéaire du 1\textsuperscript{er} ordre ?
\quad A)~$y'' + 2y = 0$
\quad B)~$y \cdot y' = x$
\quad C)~$y' + 3y = e^x$
\quad D)~$(y')^2 = y$}

\begin{methodebox}
\textbf{EDO linéaire du 1\textsuperscript{er} ordre.} Forme standard : $y' + p(x)\,y = q(x)$, où $y$ et $y'$ apparaissent à la puissance 1, sans produit entre eux.
\end{methodebox}

\begin{itemize}[leftmargin=1.5em, itemsep=2pt]
  \item \textcolor{rougeerr}{\textbf{A)}} $y''$ présent : c'est du 2\textsuperscript{e} ordre.
  \item \textcolor{rougeerr}{\textbf{B)}} $y\cdot y'$ : produit de $y$ et $y'$, donc non linéaire.
  \item \textcolor{vertok}{\textbf{C)}} $y'$ et $y$ à la puissance 1, pas de produit : \textbf{linéaire du 1\textsuperscript{er} ordre.} \textbf{Correct.}
  \item \textcolor{rougeerr}{\textbf{D)}} $(y')^2$ : non linéaire (puissance 2).
\end{itemize}

%────── Q20 ────────────────────────────────────────────────────────────────────
\qtitre{QCM}{D - $y = Ce^{4x}$}
\qenonce{La solution générale de $y' - 4y = 0$ est :
\quad A)~$y = Ce^{-4x}$
\quad B)~$y = 4x + C$
\quad C)~$y = C\cos(4x)$
\quad D)~$y = Ce^{4x}$}

\begin{methodebox}
\textbf{EDO linéaire homogène du 1\textsuperscript{er} ordre.} $y' + p(x)\,y = 0$ $\Rightarrow$ $y = C\,e^{-\int p(x)\,dx}$.

Ici : $y' - 4y = 0$ s'écrit $y' + (-4)y = 0$, donc $p(x) = -4$ :
$$y = C\,e^{-\int(-4)\,dx} = C\,e^{4x}$$
\end{methodebox}

La proposition A correspond à $y' + 4y = 0$ (signe inversé).

%────── Q21 ────────────────────────────────────────────────────────────────────
\qtitre{QCM}{C - $r^2 - 5r + 6 = 0$}
\qenonce{L'équation caractéristique associée à $y'' - 5y' + 6y = 0$ est :
\quad A)~$r^2 + 5r + 6 = 0$
\quad B)~$r - 5 + 6 = 0$
\quad C)~$r^2 - 5r + 6 = 0$
\quad D)~$r^2 - 5r = 0$}

\begin{methodebox}
\textbf{Équation caractéristique.} Pour $ay'' + by' + cy = 0$, on substitue $y'' \to r^2$, $y' \to r$, $y \to 1$ :
$$ar^2 + br + c = 0$$
Ici : $1\cdot r^2 + (-5)\cdot r + 6\cdot 1 = 0 \;\Rightarrow\; r^2 - 5r + 6 = 0$.
\end{methodebox}

La proposition A a les signes inversés. La proposition D oublie le terme $+6$.

%────── Q22 ────────────────────────────────────────────────────────────────────
\qtitre{QCM}{A - $y = e^{-x}(C_1\cos 2x + C_2\sin 2x)$}
\qenonce{L'équation $y'' + 2y' + 5y = 0$ admet les racines complexes $r = -1 \pm 2i$. Sa solution générale est :
\quad A)~$y = e^{-x}(C_1\cos 2x + C_2\sin 2x)$
\quad B)~$y = e^{x}(C_1\cos 2x + C_2\sin 2x)$
\quad C)~$y = (C_1 + C_2 x)\,e^{-x}$
\quad D)~$y = C_1 e^{-x} + C_2 e^{5x}$}

\begin{methodebox}
\textbf{Racines complexes conjuguées $r = \alpha \pm \beta i$.} La solution générale est :
$$y = e^{\alpha x}(C_1\cos\beta x + C_2\sin\beta x)$$
Ici $\alpha = -1$, $\beta = 2$ : $y = e^{-x}(C_1\cos 2x + C_2\sin 2x)$.
\end{methodebox}

La proposition B aurait $\alpha = +1$ (signe de la partie réelle inversé). La proposition C est la forme pour une \textbf{racine double réelle}.

%────── Q23 ────────────────────────────────────────────────────────────────────
\qtitre{QCM}{B - $y = (C_1 + C_2 x)\,e^{r_0 x}$}
\qenonce{Si l'équation caractéristique admet une racine double $r_0$, la solution générale est :
\quad A)~$y = C_1 e^{r_0 x} + C_2 e^{-r_0 x}$
\quad B)~$y = (C_1 + C_2 x)\,e^{r_0 x}$
\quad C)~$y = e^{r_0 x}(C_1\cos r_0 x + C_2\sin r_0 x)$
\quad D)~$y = C\,e^{r_0 x}$}

\begin{methodebox}
\textbf{Racine double $r_0$.} La solution comporte un facteur polynomial pour assurer l'indépendance linéaire des deux solutions : $y = (C_1 + C_2 x)\,e^{r_0 x}$.

La proposition A correspond au cas $\Delta > 0$ (deux racines réelles distinctes $r_0$ et $-r_0$). La proposition C est la forme pour des racines complexes. La proposition D n'a qu'une constante, donc une seule solution.
\end{methodebox}

%────── Q24 (Calcul) ──────────────────────────────────────────────────────────
\qtitre{Calcul}{$r_1 = 3$, $r_2 = -2$ \quad et \quad $y(x) = C_1 e^{3x} + C_2 e^{-2x}$}
\qenonce{Résoudre l'équation caractéristique $r^2 - r - 6 = 0$. En déduire la solution générale de $y'' - y' - 6y = 0$.}

\begin{reponsebox}
\textbf{Calcul du discriminant :}
$$\Delta = (-1)^2 - 4\times 1\times(-6) = 1 + 24 = 25$$

\textbf{Racines :}
$$r = \frac{1 \pm \sqrt{25}}{2} = \frac{1 \pm 5}{2} \qquad \Rightarrow \qquad r_1 = 3, \quad r_2 = -2$$

\textbf{$\Delta > 0$ : deux racines réelles distinctes} $\Rightarrow$ solution générale :
$$\boxed{y(x) = C_1 e^{3x} + C_2 e^{-2x}}$$
\end{reponsebox}

% ═══════════════════════════════════════════════════════════════════════════════
\section{Partie 5 - Extrema et intégrales doubles}
% ═══════════════════════════════════════════════════════════════════════════════

%────── Q25 ────────────────────────────────────────────────────────────────────
\qtitre{QCM}{D - $\nabla f(a,b) = \mathbf{0}$}
\qenonce{Un point $(a,b)$ est un point critique de $f$ si et seulement si :
\quad A)~$f(a,b) = 0$
\quad B)~La hessienne de $f$ en $(a,b)$ est nulle
\quad C)~$f$ n'est pas définie en $(a,b)$
\quad D)~$\nabla f(a,b) = \mathbf{0}$, c'est-à-dire $\dfrac{\partial f}{\partial x}(a,b) = \dfrac{\partial f}{\partial y}(a,b) = 0$}

\begin{methodebox}
\textbf{Point critique (ou stationnaire).} $(a,b)$ est un point critique si toutes les dérivées partielles premières s'annulent :
$$\nabla f(a,b) = \mathbf{0} \;\Longleftrightarrow\; f_x(a,b) = 0 \;\text{ et }\; f_y(a,b) = 0$$
C'est une \textbf{condition nécessaire} pour un extremum local. La hessienne sert ensuite à \emph{classifier} le point.
\end{methodebox}

\begin{piegebox}
\textbf{Piège fréquent.} La valeur $f(a,b)$ n'a pas à être nulle en un point critique. Exemple : $f(x,y)=x^2+y^2+5$ a un minimum en $(0,0)$ avec $f(0,0)=5\neq 0$.
\end{piegebox}

%────── Q26 ────────────────────────────────────────────────────────────────────
\qtitre{QCM}{B - Un minimum local}
\qenonce{Soit $H = \dfrac{\partial^2 f}{\partial x^2}\dfrac{\partial^2 f}{\partial y^2} - \!\left(\dfrac{\partial^2 f}{\partial x\partial y}\right)^{\!2}$ le discriminant hessien.
Si $H > 0$ et $\dfrac{\partial^2 f}{\partial x^2}(a,b) > 0$, alors $(a,b)$ est :
\quad A)~Un point selle
\quad B)~Un minimum local
\quad C)~On ne peut pas conclure
\quad D)~Un maximum local}

\begin{methodebox}
\textbf{Critère du discriminant hessien $H$.} En un point critique de $f\in\mathcal{C}^2$ :
\begin{center}
\begin{tabular}{@{}lll@{}}
\toprule
Condition & Nature & \\
\midrule
$H > 0$ et $f_{xx} > 0$ & Minimum local & paraboloïde vers le haut \\
$H > 0$ et $f_{xx} < 0$ & Maximum local & paraboloïde vers le bas \\
$H < 0$                 & Point selle   & \\
$H = 0$                 & Cas douteux   & critère insuffisant \\
\bottomrule
\end{tabular}
\end{center}
\end{methodebox}

%────── Q27 ────────────────────────────────────────────────────────────────────
\qtitre{QCM}{C - Coordonnées polaires : $x=r\cos\theta,\ y=r\sin\theta$}
\qenonce{Pour calculer $\displaystyle\iint_D \sqrt{x^2+y^2}\,dA$ sur le disque $x^2+y^2\leq 9$, le changement de variables le plus adapté est :
\quad A)~$x=u+v,\ y=u-v$
\quad B)~$x=r\cosh\theta,\ y=r\sinh\theta$
\quad C)~Coordonnées polaires : $x=r\cos\theta,\ y=r\sin\theta$
\quad D)~Aucun changement nécessaire}

\begin{methodebox}
\textbf{Coordonnées polaires.} Sur le disque $r \leq R$ avec $\sqrt{x^2+y^2} = r$ :
$$\iint_D \sqrt{x^2+y^2}\,dA = \int_0^{2\pi}\!\int_0^3 r \cdot r\,dr\,d\theta = \int_0^{2\pi}\!\int_0^3 r^2\,dr\,d\theta = 2\pi\cdot\frac{27}{3} = 18\pi$$
Le jacobien du changement de variables polaires est $\left|\dfrac{\partial(x,y)}{\partial(r,\theta)}\right| = r$, d'où $dA = r\,dr\,d\theta$.
\end{methodebox}

%────── Q28 ────────────────────────────────────────────────────────────────────
\qtitre{QCM}{D - $\displaystyle\int_0^1 x\,dx \cdot \int_0^2 y\,dy$}
\qenonce{On calcule $\displaystyle\iint_D xy\,dA$ sur $D=[0,1]\times[0,2]$. Par séparation de Fubini :
\quad A)~$\displaystyle\int_0^1\!\int_0^2 (x+y)\,dy\,dx$
\quad B)~$\displaystyle\Bigl(\int_0^1 xy\,dx\Bigr)^2$
\quad C)~$\displaystyle 2\int_0^1\!\int_0^2 xy\,dy\,dx$
\quad D)~$\displaystyle\int_0^1 x\,dx \cdot \int_0^2 y\,dy$}

\begin{methodebox}
\textbf{Théorème de Fubini (cas séparable).} Sur un rectangle $[a,b]\times[c,d]$, si $f(x,y) = g(x)\cdot h(y)$ :
$$\iint_D f\,dA = \int_a^b g(x)\,dx \cdot \int_c^d h(y)\,dy$$
Ici $xy = x \cdot y$, donc $\displaystyle\iint_D xy\,dA = \int_0^1 x\,dx \cdot \int_0^2 y\,dy = \frac{1}{2}\cdot 2 = 1$.
\end{methodebox}

%────── Q29 ────────────────────────────────────────────────────────────────────
\qtitre{QCM}{C - Un point selle}
\qenonce{Si le discriminant hessien $H < 0$ en un point critique $(a,b)$, alors $(a,b)$ est :
\quad A)~Un minimum local
\quad B)~Un maximum local
\quad C)~Un point selle
\quad D)~On ne peut pas conclure}

\begin{methodebox}
\textbf{Point selle.} Si $H < 0$, la hessienne est indéfinie (valeurs propres de signes opposés) : $f$ croît dans certaines directions et décroît dans d'autres depuis $(a,b)$. C'est un col, pas un extremum.
\end{methodebox}

\textbf{Exemple canonique :} $f(x,y) = x^2 - y^2$ en $(0,0)$ : $\nabla f(0,0) = \mathbf{0}$, $H = (2)(−2) - 0^2 = -4 < 0$ $\Rightarrow$ point selle.

%────── Q30 (Calcul) ──────────────────────────────────────────────────────────
\qtitre{Calcul}{$\displaystyle\int_0^2\!\int_0^2 (x+y)\,dy\,dx = 8$}
\qenonce{Calculer $\displaystyle\int_0^2\!\int_0^2 (x+y)\,dy\,dx$.}

\begin{reponsebox}
\textbf{Étape 1 - intégration sur $y$ (intégrale intérieure) :}
$$\int_0^2(x+y)\,dy = \Bigl[xy+\frac{y^2}{2}\Bigr]_0^2 = \bigl(2x + 2\bigr) - 0 = 2x + 2$$

\textbf{Étape 2 - intégration sur $x$ (intégrale extérieure) :}
$$\int_0^2(2x+2)\,dx = \Bigl[x^2+2x\Bigr]_0^2 = (4 + 4) - 0 = \boxed{8}$$
\end{reponsebox}

% ═══════════════════════════════════════════════════════════════════════════════
\newpage
\section*{\color{bleutexte}Synthèse des méthodes à retenir}
% ═══════════════════════════════════════════════════════════════════════════════

\begin{methodebox}
\textbf{Fonctions de plusieurs variables - Récapitulatif}
\begin{enumerate}[itemsep=2pt, topsep=2pt]
  \item Continuité en $a$ : $\lim_{(x,y)\to a}f(x,y) = f(a)$. Différentiable $\Rightarrow$ continu (réciproque fausse).
  \item Dérivée partielle $\partial_x f$ : dériver $f$ en ne faisant varier que $x$, $y$ fixé.
  \item Différentielle : $df = \dfrac{\partial f}{\partial x}dx + \dfrac{\partial f}{\partial y}dy$.
  \item Plan tangent à $z=f(x,y)$ en $(a,b)$ : $z = f(a,b) + f_x(a,b)(x-a) + f_y(a,b)(y-b)$.
  \item Règle de la chaîne : $\dfrac{\partial(g\circ h)}{\partial x} = g'(h)\cdot\dfrac{\partial h}{\partial x}$.
  \item Schwartz ($f\in\mathcal{C}^2$) : $\dfrac{\partial^2 f}{\partial x\partial y} = \dfrac{\partial^2 f}{\partial y\partial x}$.
\end{enumerate}
\end{methodebox}

\vspace{6pt}

\begin{methodebox}
\textbf{Opérateurs vectoriels - Récapitulatif}
\begin{center}
\begin{tabular}{@{}llll@{}}
\toprule
Opérateur & Notation & Formule & Nature \\
\midrule
Gradient   & $\nabla f$               & $\bigl(\partial_x f,\,\partial_y f,\,\partial_z f\bigr)$                                         & Vecteur \\[4pt]
Divergence & $\text{div}\,\mathbf{F}$ & $\partial_x P + \partial_y Q + \partial_z R$                                                    & Scalaire \\[4pt]
Rotationnel & $\text{rot}\,\mathbf{F}$ & $\bigl(\partial_y R - \partial_z Q,\;\partial_z P - \partial_x R,\;\partial_x Q - \partial_y P\bigr)$ & Vecteur \\[4pt]
Laplacien  & $\Delta f$               & $\partial^2_{xx} f + \partial^2_{yy} f + \partial^2_{zz} f$                                    & Scalaire \\
\bottomrule
\end{tabular}
\end{center}
\textbf{Identités :} $\text{rot}(\nabla f) = \mathbf{0}$ \quad et \quad $\text{div}(\text{rot}\,\mathbf{F}) = 0$
\end{methodebox}

\vspace{6pt}

\begin{methodebox}
\textbf{Équations différentielles - Récapitulatif}
\begin{center}
\begin{tabular}{@{}lll@{}}
\toprule
Type de racines ($r^2+br+c=0$) & Condition & Solution générale \\
\midrule
Deux réelles distinctes $r_1 \neq r_2$ & $\Delta > 0$ & $C_1 e^{r_1 x} + C_2 e^{r_2 x}$ \\[4pt]
Racine double $r_0$              & $\Delta = 0$ & $(C_1 + C_2 x)\,e^{r_0 x}$ \\[4pt]
Complexes $\alpha \pm \beta i$   & $\Delta < 0$ & $e^{\alpha x}(C_1\cos\beta x + C_2\sin\beta x)$ \\
\bottomrule
\end{tabular}
\end{center}
\end{methodebox}

\vspace{6pt}

\begin{methodebox}
\textbf{Extrema - Classification des points critiques}

En un point critique ($\nabla f = \mathbf{0}$), calculer $H = f_{xx}\,f_{yy} - f_{xy}^2$ :
\begin{center}
\begin{tabular}{@{}lll@{}}
\toprule
Condition & Nature & \\
\midrule
$H > 0$ et $f_{xx} > 0$ & Minimum local    & \\
$H > 0$ et $f_{xx} < 0$ & Maximum local    & \\
$H < 0$                 & Point selle (col) & \\
$H = 0$                 & Cas douteux       & critère insuffisant \\
\bottomrule
\end{tabular}
\end{center}
\end{methodebox}

\end{document}
